\section{提示框与强调}
\label{sec:00-boxes}

\subsection{要点框}

需要把一小段话单独拎出来提醒读者时，用 \verb|\KeyPoint{标题}{说明}|。
它不加框线，靠颜色和字号区分，不会打断阅读节奏。

写法：

\begin{lstlisting}[language=TeX]
\KeyPoint{安全边界}{没有把握的结论不要写成事实，
用"待补充"占位并注明来源，等实验数据出来再补。}
\end{lstlisting}

效果：

\KeyPoint{安全边界}{没有把握的结论不要写成事实，用"待补充"占位并注明来源，
等实验数据出来再补。}

常见用法有两种：一节开头用它点出"这一节解决什么问题"，或者一段关键结论后面
用它补一句提醒。\textbf{一页里最多一两个}，多了就不"要点"了。

\subsection{其他强调手段}

\begin{description}
  \item[加粗] \verb|\textbf{...}|，用于术语第一次出现和关键结论
  \item[强调色] \verb|\textcolor{Accent}{...}|，用于需要视觉引导的关键词
  \item[缩小字号] \verb|{\small ...}|，用于附注性的补充说明
  \item[引用块] \texttt{quote} 环境，用于成段引文或需要独立成段的话
\end{description}

\subsection{一节的实验怎么排}

仓库目前没有专门的"实验"环境，约定用固定四段式：目标、环境、步骤、验收。
步骤用有序列表，验收用可勾选清单：

\begin{lstlisting}[language=TeX]
\textbf{实验目标}\quad 在 FPGA 上点亮流水线各级的指示灯。

\textbf{实验环境}\quad Vivado 2024.1；开发板型号见 config/hardware.tex。

\textbf{实验步骤}
\begin{enumerate}
  \item 综合并生成比特流；
  \item 下载到开发板；
  \item 观察指示灯的流水效果。
\end{enumerate}

\textbf{验收清单}
\begin{itemize}
  \item[$\square$] 综合没有关键警告
  \item[$\boxtimes$] 指示灯按预期流水点亮
\end{itemize}
\end{lstlisting}

如果你觉得这套四段式应该固化成统一环境，\textbf{在 Issue 里提出来}，
由主编加到 \texttt{styles/} 里，不要自己在章节里拼一个。
